\subsection{Математичне обґрунтування та виведення барицентричної інтерполяційної формули}

Фундаментальне вивчення поліноміальної інтерполяції має глибоке історичне коріння. Класичним підходом до розв'язання цієї задачі є використання альтернативної форми Лагранжа, де інтерполяційний поліном записується як лінійна комбінація базисних (фундаментальних) поліномів Лагранжа:
\begin{equation}\label{eq:lagrange_form}
	p(x) = \sum_{j=0}^{n} f_j \ell_j(x).
\end{equation}

Нехай задано набір різних інтерполяційних вузлів $x_0, \dots, x_n$, які можуть бути як дійсними, так і комплексними. Тоді $j$-й поліном Лагранжа $\ell_j(x)$ є єдиним поліномом ступеня $n$, який набуває значення $1$ у вузлі $x_j$ та дорівнює $0$ у всіх інших вузлах $x_k$:
\begin{equation}\label{eq:lagrange_basis_def}
	\ell_j(x_k) = 
	\begin{cases} 
		1, & k = j, \\ 
		0, & k \neq j. 
	\end{cases}
\end{equation}

Аналітичний вираз для базисного полінома $\ell_j(x)$ можна легко записати у явному вигляді:
\begin{equation}\label{eq:lagrange_basis_explicit}
	\ell_j(x) = \frac{\prod_{k \neq j} (x - x_k)}{\prod_{k \neq j} (x_j - x_k)}.
\end{equation}

Оскільки знаменник у виразі~\eqref{eq:lagrange_basis_explicit} є константою, ця функція дійсно є поліномом ступеня $n$ із заданими нулями, і очевидно, що вона дорівнює одиниці при $x = x_j$. Рівняння~\eqref{eq:lagrange_basis_explicit} є стандартним і загальновідомим представленням інтерполяції Лагранжа, яке широко зустрічається в академічній літературі.

З обчислювальної точки зору формула~\eqref{eq:lagrange_form} є математично коректною, проте вираз~\eqref{eq:lagrange_basis_explicit} має суттєві недоліки при програмній реалізації. Розрахунок $\ell_j(x)$ для кожного значення $x$ вимагає виконання $O(n)$ операцій. Оскільки у формулі~\eqref{eq:lagrange_form} необхідно виконати $O(n)$ таких обчислень та додати їх результати, загальна алгоритмічна складність знаходження значення $p(x)$ в одній точці $x$ становить $O(n^2)$.

Шляхом алгебраїчних перетворень можна суттєво оптимізувати кількість необхідних обчислювальних операцій. Ключове спостереження полягає в тому, що для різних значень індексу $j$ чисельники у виразі~\eqref{eq:lagrange_basis_explicit} є ідентичними, за винятком відсутності відповідного множника $(x - x_j)$. Щоб використати цю спільність, введемо так званий вузловий поліном $\ell(x)$ ступеня $n+1$ для заданої сітки:
\begin{equation}\label{eq:nodal_polynomial}
	\ell(x) = \prod_{k=0}^{n} (x - x_k).
\end{equation}

Тоді вираз~\eqref{eq:lagrange_basis_explicit} можна переписати у вигляді елементарної, але надзвичайно важливої тотожності:
\begin{equation}\label{eq:lagrange_identity}
	\ell_j(x) = \frac{\ell(x)}{\ell'(x_j)(x - x_j)}.
\end{equation}

Введемо додаткове позначення для вагових коефіцієнтів:
\begin{equation}\label{eq:barycentric_weights_def}
	\lambda_j = \frac{1}{\prod_{k \neq j} (x_j - x_k)},
\end{equation}
або, що є еквівалентним:
\begin{equation}\label{eq:barycentric_weights_derivative}
	\lambda_j = \frac{1}{\ell'(x_j)}.
\end{equation}

З урахуванням введених вагових коефіцієнтів, рівняння~\eqref{eq:lagrange_basis_explicit} набуває вигляду:
\begin{equation}\label{eq:lagrange_basis_barycentric}
	\ell_j(x) = \ell(x) \frac{\lambda_j}{x - x_j},
\end{equation}

а класична формула Лагранжа~\eqref{eq:lagrange_form} трансформується у наступний вираз:
\begin{equation}\label{eq:barycentric_first_form}
	p(x) = \ell(x) \sum_{j=0}^{n} \frac{\lambda_j}{x - x_j} f_j.
\end{equation}

Рівняння~\eqref{eq:barycentric_first_form} в сучасній літературі називається «модифікованою формулою Лагранжа» або \textbf{«першою формою барицентричної інтерполяційної формули»}. Її головна цінність полягає в тому, що залежність від координати $x$ всередині суми є максимально простою. Якщо вагові коефіцієнти $\{\lambda_j\}$ заздалегідь відомі, розрахунок значення $p(x)$ за формулою~\eqref{eq:barycentric_first_form} вимагає лише $O(n)$ операцій. Розрахунок самих ваг за формулою~\eqref{eq:barycentric_weights_def} потребує $O(n^2)$ операцій, проте ця процедура виконується лише один раз для заданої сітки вузлів $\{x_j\}$ незалежно від значень $x$. Більше того, для спеціальних сіток (таких як вузли Чебишева), ваги відомі аналітично і їх взагалі не потрібно обчислювати ітеративно.

Незважаючи на ефективність першої форми, існує ще більш елегантна математична модель. Відомо, що сума всіх базисних поліномів Лагранжа $\ell_j$ утворює поліном ступеня $n$, який набуває значення $1$ у кожному вузлі сітки. Оскільки поліноміальна інтерполяція є єдиною, цією сумою може бути лише поліном-константа, що тотожно дорівнює одиниці:
\begin{equation}\label{eq:lagrange_sum_one}
	\sum_{j=0}^{n} \ell_j(x) = 1.
\end{equation}

Якщо поділити модифікований вираз~\eqref{eq:lagrange_basis_barycentric} на цю тотожність, спільний множник $\ell(x)$ математично скорочується, що дає змогу вивести фінальну залежність:
\begin{equation}\label{eq:barycentric_second_form_basis}
	\ell_j(x) = \frac{\frac{\lambda_j}{x - x_j}}{\sum_{k=0}^{n} \frac{\lambda_k}{x - x_k}}.
\end{equation}

Підставивши це представлення назад у вихідне рівняння~\eqref{eq:lagrange_form}, ми отримуємо класичну \textbf{«другу форму барицентричної інтерполяційної формули»} (або «справжню барицентричну формулу») для поліноміальної інтерполяції на довільному наборі з $n+1$ вузлів $\{x_j\}$.

З аналізу виведеної другої барицентричної формули~\eqref{eq:barycentric_second_form_basis} випливає, що вона коректно інтерполює вихідні дані. Якщо координата $x$ наближається до одного з вузлів $x_j$, відповідні доданки в чисельнику та знаменнику прямують до нескінченності, проте їхнє відношення в границі дає точне значення $f_j$. Неочевидним, але фундаментальним є той факт, що ця функція, маючи зовнішній вигляд раціонального дробу, насправді є поліномом ступеня $n$. Ця властивість гарантується виключно завдяки спеціальним значенням вагових коефіцієнтів $\lambda_j$. 

Для екстремумів поліномів Чебишева (вузлів Чебишева-Лобатто) ці вагові коефіцієнти набувають доволі простого аналітичного вигляду: вони дорівнюють $(-1)^j$, помноженому на загальну константу $2^{n-1}/n$, причому для крайніх вузлів сітки ($j = 0$ та $j = n$) це значення додатково зменшується вдвічі.

%
У результаті спрощення вагових коефіцієнтів ми отримуємо надзвичайно просту та обчислювально ефективну барицентричну формулу інтерполяції за вузлами Чебишева:
\begin{equation}\label{eq:cheb_barycentric}
	p(x) = \frac{\sum\limits_{j=0}^{n} \vphantom{\sum}^{\prime} \dfrac{(-1)^j}{x - x_j} f_j}{\sum\limits_{j=0}^{n} \vphantom{\sum}^{\prime} \dfrac{(-1)^j}{x - x_j}},
\end{equation}
де штрих біля знака суми ($\sum^{\prime}$) є традиційним математичним позначенням, яке вказує на те, що перший ($j=0$) та останній ($j=n$) доданки ряду додатково множаться на множник $1/2$. Для особливого випадку, коли $x = x_j$, значення полінома тотожно приймається рівним $f_j$.

Представлене рівняння~\eqref{eq:cheb_barycentric} має важливу властивість масштабної інваріантності (scale-invariance): при масштабуванні сітки вузлів Чебишева на будь-який довільний інтервал $[a, b]$, математичний вигляд формули залишається абсолютно незмінним. З точки зору комп'ютерної реалізації алгоритму це є вагомою перевагою, оскільки гарантує відсутність проблем переповнення (overflow) або втрати значущості (underflow) розрядної сітки процесора, навіть якщо $n$ вимірюється тисячами.

\begin{proof}
	Рівняння~\eqref{eq:cheb_barycentric} є окремим випадком загальної барицентричної формули. Для його доведення покажемо, що для вузлів Чебишева вагові коефіцієнти $\lambda_j$ зводяться до константи $\frac{2^{n-1}}{n}$, помноженої на $(-1)^j$, та половини цього значення для граничних точок $j=0$ або $n$.
	
	Почнемо з того, що для вузлів Чебишева вузловий поліном $\ell(x)$ можна виразити через поліноми Чебишева першого роду $T_k(x)$ як:
	\begin{equation}
		\ell(x) = 2^{-n} (T_{n+1}(x) - T_{n-1}(x)).
	\end{equation}
	З урахуванням визначення базисних поліномів, це означає:
	\begin{equation}
		\ell_j(x) = 2^{-n} \lambda_j \frac{T_{n+1}(x) - T_{n-1}(x)}{x - x_j}.
	\end{equation}
	Оскільки барицентричні ваги визначаються як зворотна величина до похідної вузлового полінома $\lambda_j = \frac{1}{\ell'(x_j)}$, маємо:
	\begin{equation}
		\lambda_j = \frac{1}{\ell'(x_j)} = \frac{2^n}{T'_{n+1}(x_j) - T'_{n-1}(x_j)}.
	\end{equation}
	Використовуючи властивості похідних поліномів Чебишева, можна показати, що у вузлах екстремумів різниця похідних становить:
	\begin{equation}
		T'_{n+1}(x_j) - T'_{n-1}(x_j) = 2n(-1)^j, \quad 1 \le j \le n-1,
	\end{equation}
	причому для граничних точок $j=0$ та $n$ це значення подвоюється. Таким чином, отримуємо вагові коефіцієнти:
	\begin{equation}\label{eq:lambda_cheb_proof}
		\lambda_j = \frac{2^{n-1}}{n}(-1)^j, \quad 1 \le j \le n-1,
	\end{equation}
	і половину цього значення для $j=0$ та $n$.
	
	Оскільки у другій барицентричній формулі коефіцієнти $\lambda_j$ фігурують одночасно і в чисельнику, і в знаменнику, загальна константа $\frac{2^{n-1}}{n}$ математично скорочується. У результаті залишається лише множник $(-1)^j$, що й вимагалося довести.
\end{proof}
%
\subsection{Аналіз числової стійкості барицентричної інтерполяції}

Експериментальні та теоретичні дослідження доводять, що барицентрична інтерполяція за вузлами Чебишева є надзвичайно надійним обчислювальним процесом. Незважаючи на побоювання щодо можливих помилок округлення (cancellation errors) машинної арифметики у випадках, коли знаменники внутрішніх дробів $(x - x_j)$ прямують до нуля при $x \to x_j$, алгоритм довів свою абсолютну числову стійкість у форматі з плаваючою комою для всіх $x \in [-1, 1]$. Це кардинально відрізняє його від класичного алгоритму поліноміальної інтерполяції через розв'язання лінійної системи рівнянь Вандермонда, який є експоненціально нестабільним.

Відповідно до фундаментальних досліджень з чисельного аналізу, при проєктуванні обчислювальних алгоритмів важливо розмежовувати властивості першої та другої барицентричних форм. Перша форма барицентричної формули володіє сильною властивістю зворотної стійкості (backward stability) і залишається стабільною навіть при екстраполяції даних (за межами інтервалу $[-1, 1]$). Однак її критичним практичним недоліком є відсутність масштабної інваріантності: вагові коефіцієнти зростають експоненціально (пропорційно $2^n$). У стандартній комп'ютерній арифметиці подвійної точності (IEEE double precision) це неминуче призводить до переповнення розрядної сітки (overflow), якщо кількість інтерполяційних вузлів $n$ перевищує 1000.

Натомість друга барицентрична формула для вузлів Чебишева~\eqref{eq:cheb_barycentric} повністю позбавлена проблеми алгоритмічного переповнення завдяки скороченню загального множника в чисельнику та знаменнику. Вона математично гарантує пряму стійкість (forward stability) за умови, що оцінка полінома здійснюється всередині базового інтервалу $[-1, 1]$, а точки дискретизації розподілені за законом Чебишева. 

Таким чином, для більшості прикладних задач комп'ютерної графіки, зокрема для просторового масштабування цифрових зображень, друга барицентрична формула є оптимальним та безальтернативним вибором. Вона дозволяє максимально стабільно обчислювати значення інтерполяційного полінома навіть для растрових матриць, що містять тисячі або мільйони вузлів. 